\DocumentMetadata{
  pdfversion=2.0,
  pdfstandard=ua-2,
  lang=en-US,
  tagging=on,
  tagging-setup={math/setup=mathml-SE,tabsorder=structure}
}
\documentclass[11pt,letterpaper]{article}
\usepackage[margin=0.68in,includefoot]{geometry}
\usepackage{newcomputermodern}
\usepackage{amsmath,amscd,mathtools}
\providecommand{\square}{\mdlgwhtsquare}
\usepackage[hidelinks]{hyperref}
\hypersetup{
  pdftitle={Algebra PhD Exam, January 2003},
  pdfauthor={University of Florida Department of Mathematics},
  pdfsubject={Graduate mathematics examination},
  pdfdisplaydoctitle=true,
  bookmarksopen=true
}
\setcounter{secnumdepth}{0}
\setlength{\parindent}{0pt}
\setlength{\parskip}{0.28em}
\setlength{\emergencystretch}{3em}
\setlength{\leftmargini}{2.15em}
\setlength{\leftmarginii}{2.65em}
\setlength{\leftmarginiii}{3em}
\renewcommand{\labelenumi}{\textbf{\theenumi.}}
\pagestyle{empty}
\raggedbottom
\allowdisplaybreaks
\newenvironment{examproblems}[1]{%
  \begin{enumerate}%
  \setcounter{enumi}{#1}%
  \setlength{\topsep}{0.35em}%
  \setlength{\partopsep}{0pt}%
  \setlength{\itemsep}{0.48em}%
  \setlength{\parsep}{0.18em}%
}{\end{enumerate}}
\newenvironment{examsubparts}{%
  \begin{enumerate}%
  \setlength{\topsep}{0.18em}%
  \setlength{\partopsep}{0pt}%
  \setlength{\itemsep}{0.16em}%
  \setlength{\parsep}{0.1em}%
}{\end{enumerate}}
\newenvironment{examinstructions}{%
  \par\begingroup\itshape
}{%
  \par\endgroup\vspace{0.18em}
}
\makeatletter
\renewcommand\section{\@startsection{section}{1}{0pt}%
  {-1.25ex plus -.25ex minus -.1ex}%
  {0.55ex plus .1ex}%
  {\normalfont\large\bfseries}}
\renewcommand{\@maketitle}{%
  \newpage\null\vskip -1.2em%
  {\footnotesize\bfseries
    \href{https://www.ufl.edu/}{University of Florida}, \href{https://math.ufl.edu/}{Department of Mathematics}\par}%
  \vskip 0.18em%
  \tagstructbegin{tag=Title}%
    {\LARGE\bfseries\href{https://gma.math.ufl.edu/exams/algebra-phd/}{Algebra PhD Exam}, January 2003\par}%
  \tagstructend%
  \vskip 0.35em%
  \tagmcbegin{artifact}%
    \hrule height 1pt%
  \tagmcend%
  \vskip 0.55em%
}
\makeatother
\title{Algebra PhD Exam, January 2003}
\author{}
\date{}
\begin{document}
\maketitle
\thispagestyle{empty}
\begin{examinstructions}
Do seven of the following eleven problems. Please do not turn in more than seven problems. You must show your work. Answers with no work and/or no explanations will receive no credit. State clearly any theorem you use in your proofs. In the problems, \(\mathbb Z\), resp. \(\mathbb Q\), \(\mathbb C\), is the set of all integers, resp. of all rational numbers, of all complex numbers.
\end{examinstructions}
\begin{examproblems}{0}
\item
State and prove the Jordan–Hölder Theorem for finite groups.
\item
Let~\(F\) be the free group in two generators~\(a\) and~\(b\). Let~\(N\) be the smallest normal subgroup of~\(F\) which contains \(a^5\), \(b^2\) and \(abab\). Prove that \(F/N\) is a non-abelian group of order 10.
\item
Let~\(G\) be an abelian group. Define the torsion subgroup \(G_t\) of~\(G\). Prove from your definition that \(G_t\) is a subgroup of~\(G\) and that \(G/G_t\) is torsion free.
\item
Let~\(\zeta\) be a primitive 2003-rd root of unity. Prove that \(\mathbb Q(\zeta)\) is a Galois extension of \(\mathbb Q\). Calculate the Galois group of \(\mathbb Q(\zeta)/\mathbb Q\) and the degree \([\mathbb Q(\zeta):\mathbb Q]\) of the extension. Furthermore, calculate the degree \([\mathbb Q(\zeta)\cap\mathbb R:\mathbb Q]\). (Hint: 2003 is a prime).
\item
A commutative ring~\(R\) with identity is said to be Noetherian if every ideal of~\(R\) is finitely generated. Let~\(R\) be a Noetherian ring, and let~\(C\) be an ascending chain of ideals of~\(R\). Prove that~\(C\) only contains a finite number of distinct ideals.
\item
Prove the following version of the Lying-Over Theorem. Let~\(S\) be an integral extension ring of a commutative ring~\(R\) with identity. Let~\(P\) be a prime ideal of~\(R\). Prove that there exists some prime ideal~\(Q\) of~\(S\) such that \(Q\cap R=P\).
\item
Let \(A=\mathbb Z/2\mathbb Z\) and \(B=\mathbb Z/3\mathbb Z\) be viewed as \(\mathbb Z\)-modules. Calculate the \(\mathbb Z\)-module \(A\otimes B\), and, in particular, give the cardinality of \(A\otimes B\).
\item
Among the following integral domains, decide which ones are Dedekind domains, and give a brief explanation.
\begin{examsubparts}
\item[\textnormal{(a)}]
\(\mathbb Z[T]\), the polynomial ring over the integers, in one variable.
\item[\textnormal{(b)}]
\(\mathbb Z[\sqrt{-5}]=\{a+b\sqrt{-5}:a,b\in\mathbb Z\}\).
\item[\textnormal{(c)}]
The ring \(k[[T]]\) of all formal power series in one variable, over the field~\(k\).
\item[\textnormal{(d)}]
\(k[T_1,T_2]\), the polynomial ring in two variables, over the field~\(k\).
\end{examsubparts}
\item
Give one representative for each conjugacy class of elements of order 4 in \(GL(4,\mathbb Q)\).
\item
Calculate the number of irreducible polynomials of degree 4 over the field of 2 elements.
\item
Let~\(M\) be a semi-simple module over a ring~\(R\). Prove that every homomorphic image of~\(M\) is semi-simple.
\end{examproblems}
\end{document}
